\section{Fazit}

\subsection{Zusammenfassung}

Im Rahmen dieses Projekts wurde ein vollständig funktionsfähiger
\textbf{Dispatcher} für das verteilte Task-Processing-System TaskGrid+
entwickelt und containerisiert.

Die folgenden Kernkomponenten wurden implementiert:

\begin{itemize}
  \item \textbf{Zustandsmaschine} mit 8 Zuständen und vollständiger
        \texttt{InvalidTransitionError}-Absicherung.
  \item \textbf{Dispatch-Loop} als Background-Thread mit ThreadPoolExecutor
        für parallele Verarbeitung.
  \item \textbf{Timeout- und Retry-Mechanismus} mit konfigurierbaren
        Parametern und vollständiger Zustandsübergang-Unterstützung.
  \item \textbf{Namensdienst-Integration}: dynamische Worker-Adressauflösung,
        keine statisch hinterlegten Adressen.
  \item \textbf{Round-Robin Worker-Selektion} für gleichmäßige Lastverteilung.
  \item \textbf{Strukturiertes Logging} gemäß §13 der Aufgabenstellung.
  \item \textbf{HTTP Monitoring-Endpoint} (GET /status, Port 8080).
  \item \textbf{Thread-sichere Datenhaltung} (TaskStore, TaskQueue).
  \item \textbf{Vollständige Containerisierung}: vier Dockerfiles,
        docker-compose.yml mit Health Checks und Skalierungsunterstützung.
\end{itemize}

\subsection{Ergebnisse}

\textbf{Phase 2 (Dispatcher)}: 49/49 Integrationstests bestanden. \\
\textbf{Phase 3 (Docker)}: Alle Dockerfiles und docker-compose.yml
vollständig implementiert und dokumentiert.

\subsection{Ausblick}

Nach Fertigstellung aller Komponenten kann das System als Ganzes
mit \texttt{docker compose up} gestartet werden.
Der Demo-Client sendet automatisch Tasks aller Pflichttypen,
die über Worker verarbeitet und zurückgeliefert werden.
Die Skalierung einzelner Worker-Typen (\texttt{--scale worker-sum=3})
ist durch das \texttt{entrypoint.sh}-basierte
\path{WORKER_ID}- Generierungsschema vollständig vorbereitet.
